% HD-F1-I05. Root complete-proof draft; not yet frozen or promoted.
\section[HD-F1-I05: exact unit dichotomy and four more rational strata]{HD-F1-I05: an exact polynomial test for every rational nine-letter four-value stratum with four to eight reflections}
\label{sec:f1-nine-letter-unit-dichotomy}

Retain the original faithful monomial tower, primitive character,
polarization and actual period comparison of HD-I01--HD-I21.
The parameters here are
\[
 n=5,\qquad m=9,\qquad g(Y)=0,\qquad \nu=4,
 \qquad A=b-\nu\in\{4,5,6,7,8\}.
\]
Write $B$ for the four nonreflection values and $Q$ for
the reduced divisor of their uncolored cycles. Unramified
uncolored cycles are included. Its degree is $8+A$.
Choose homogeneous coordinates on the actual source and
target and write the actual map as $h=[p_0:p_1]$, with
coprime degree-nine polynomials. An isomorphism
$\mathcal O_Y(Q)\simeq\mathcal O_Y(8+A)$ carries the
literal unit to a nonzero polynomial $q_Q$ with divisor $Q$.
For $V_d=H^0(\mathbf P^1,\mathcal O(d))$, define the actual
linear map and vector
\[
 M_h:V_{A-1}\oplus V_{A-1}\longrightarrow V_{8+A},
       \quad(\alpha,\beta)\longmapsto\alpha p_0+\beta p_1,
       \qquad q_Q\in V_{8+A}.
 \tag{F1I05.1}
\]
It is injective. Indeed a relation gives $p_0\mid\beta$
and $p_1\mid\alpha$, while $A-1<9$.
The exact original period rank is
\[
\begin{array}{c|c|c}
 \operatorname{rank}[M_h\mid q_Q]&
       \operatorname{rank}d\mathcal F&\operatorname{rank}d\operatorname{Per}\\ \hline
 2A+1&A+1&A+1\\
 2A&A&A
\end{array}
 \tag{F1I05.2}
\]
These are ranks of explicit coefficient matrices of homogeneous
polynomials. Their definition retains the prescribed unit and the
original map. The second row has the further necessary restriction
\[
             |Q\cap h^{-1}(b_i)|\le A-1
                       \qquad\hbox{for every }b_i\in B.
 \tag{F1I05.3}
\]
Consequently every actual ramification profile with one of these
counts at least $A$ has original period rank $A+1$ at every
smooth point. Four full $G(3,1,9)$ tuples below have respectively
$A=4,5,6,7$ and such counts. They establish nonempty strata of
exact period ranks five, six, seven and eight. No existence is
asserted for the second row of (F1I05.2).

\paragraph{1. Evaluation and the actual prescribed-unit quotient.}
Put $\mathcal V=h_*\mathcal O_Y(Q)$.
The exact finite-flat and fibre-evaluation argument in I03 applies
throughout $0\le A\le8$, as the relevant degrees are still
$A-1\ge-1$ and $A-10<0$. Explicitly, restriction to every
length-nine scheme fibre is onto since its preceding line bundle
has degree $A-1$. Thus $\mathcal V$ is generated by its
$9+A$ global sections. The evaluation morphism
\[
 e_+:H^0(\mathcal V(-1))\otimes\mathcal O(1)
                 \longrightarrow\mathcal V
\]
has fibre kernel $H^0(\mathcal O_Y(Q)\otimes h^*\mathcal O(-2))=0$.
Its source has rank $A$. Local invertible minors identify its
image with a direct summand, and its cokernel is a generated
rank-$9-A$ bundle of degree zero. Its section space has exactly
$9-A$ dimensions, by the section sequence and
$H^1(\mathcal O(1))=0$. Evaluation therefore identifies that
cokernel with $\mathcal O^{9-A}$. Lifting its basis gives
the actual isomorphism
\[
 \mathcal V\simeq\mathcal O(1)^A\oplus\mathcal O^{9-A},
 \qquad \operatorname{im}e_+=\mathcal O(1)^A.
 \tag{F1I05.4}
\]
The subbundle on the right is the original evaluation image.

The literal unit section has no zero in any fibre of $\mathcal V$.
Away from $B$ it is a unit in every local fibre algebra. Over
$b_i\in B$, a local branch has equation $w=z^\ell$. At a
point of $Q$, its coefficient relative to the generator $z^{-1}$
is $z$, nonzero modulo $z^\ell$ exactly when $\ell>1$.
At a point outside $Q$, its coefficient is $1$, always nonzero.
Vanishing of the entire fibre section would therefore require
nine distinct unramified uncolored points, making that actual
monomial inertia the identity. The branch inertia is nonidentity.
This proves the nowhere-zero assertion and recovers the precise
I21 locally free quotient $\mathcal W_Q=\mathcal V/\mathcal O$.

Write the prescribed unit in (F1I05.4) as $(s,c)$, with
$c\in\mathbf C^{9-A}$. For $c\ne0$, a constant invertible
matrix sends it to the first constant basis vector. The actual
triangular automorphism $(u,w)\mapsto(u-sw_1,w)$ then gives
\[
 \mathcal W_Q\simeq\mathcal O(1)^A\oplus\mathcal O^{8-A}.
 \tag{F1I05.5}
\]
For $c=0$, the coordinates of $s$ span all of
$H^0(\mathcal O(1))$. A smaller span would be generated by one
linear form with a zero, contradicting the established fibre
nonvanishing. A constant change of basis in $\mathbf C^A$
sends them to $(x_0,x_1,0,\ldots,0)$ for a target coordinate
basis. The literal surjection
\[
 \mathcal O(1)^2\longrightarrow\mathcal O(2),
             \qquad(u,v)\longmapsto x_0v-x_1u
\]
has kernel the actual section $(x_0,x_1)$, as is checked on
the opens $x_0\ne0$ and $x_1\ne0$. Thus this second case gives
\[
 \mathcal W_Q\simeq\mathcal O(2)\oplus
           \mathcal O(1)^{A-2}\oplus\mathcal O^{9-A}.
 \tag{F1I05.6}
\]
Every coordinate change above is an explicit automorphism of
the given bundle, with the specified unit followed through it.

\paragraph{2. Exact polynomial and original period maps.}
Under the chosen line-bundle isomorphism, the section map of
$e_+$ is precisely (F1I05.1), by pullback of target coordinates.
Therefore $c=0$ is equivalent to $q_Q\in\operatorname{im}M_h$.
Injectivity of $M_h$ gives exactly its two augmented matrix ranks.
The original I21 tangent-fibre isomorphism is
\[
              \ker d\mathcal F=H^0(\mathcal W_Q(-2)).
\]
Equations (F1I05.5)--(F1I05.6) give dimensions zero and one,
respectively. The original Hurwitz tangent has dimension
$b-3=A+1$, proving the middle column of (F1I05.2).
On the actual rational source both primitive eigendifferential
bundles have degree four. Their multiplication onto
$H^0(\mathcal O_Y(8))$ is surjective: each degree-eight monomial
is a product of two degree-four monomials. Composition with the
original I20 quadratic trace and the actual polarized isogeny
of I02 identifies its image with the original period
codifferential. The latter isogeny is the restriction of $q^*$,
with $q_*q^*=[|\ker\chi|]$ and its stated theta multiplier.
This proves the last column through the original maps.

\paragraph{3. A numerical restriction at every nonreflection value.}
In the second matrix-rank row, write
$q_Q=\alpha p_0+\beta p_1$ and take the homogeneous greatest
common factor $g$ of $\alpha,\beta$, of degree $c_0$.
The remaining coprime pair gives the actual morphism
$\phi=[-\beta/g:\alpha/g]$ of degree $e=A-1-c_0$.
For a nonconstant $\phi$, each point of $Q$ outside the
zero divisor of $g$ satisfies $h=\phi$. Over any single
$b_i\in B$, there are at most $e$ such distinct points,
since they belong to $\phi^{-1}(b_i)$. The removed divisor
contributes at most $c_0$ more distinct points. Hence their
total is at most $e+c_0=A-1$, proving (F1I05.3).

For a constant $\phi=b_0$, every point of $Q$ outside
$\operatorname{div}(g)$ lies over $b_0$. For $b_0\notin B$
this contradicts $\deg Q>c_0$. For $b_0\in B$, that fibre
has at most eight uncolored cycles: nine would again give
identity inertia. Consequently
$8+A=\deg Q\le8+c_0\le8+(A-1)$, also impossible.
This includes a zero member of the pair $\alpha,\beta$.

\paragraph{4. Four actual full-group tuples.}
Use the actual monomial matrices $C,D$ of I01. The
permutation of $C$ is $(1\,2\,3\,4\,5\,6)$, with
zero color on that cycle and color one on the fixed
letters $7,8,9$. The permutation of $D$ is $(6\,7\,8\,9)$,
of cycle color one, with fixed-letter colors one at $1$,
two at $2$, and zero at $3,4,5$.
Let $R_i$ be the diagonal matrix with entry $\zeta_3$ at
letter $i$ and one elsewhere. Let $P_j$ be the pure
permutation $(1\,\cdots\,j)$ and $t_j$ the pure
transposition $(j\,j+1)$.
With the convention that the rightmost permutation acts first,
literal multiplication gives
\[
 P_6=P_jt_jt_{j+1}\cdots t_5,\qquad
 C=P_6R_7R_8R_9,\qquad j=3,4,5.
 \tag{F1I05.7}
\]
For $A=4,5,6$, put $j=9-A$ and take the actual tuple
\[
 (P_j,t_j,t_{j+1},\ldots,t_5,
          R_7,R_8,R_9,C^{-1},D,D^{-1}).
 \tag{F1I05.8}
\]
Its initial block multiplies to $C$, so its full product is one.
All the $t_i$ and $R_i$ are reflections. Its other four
entries have ranks respectively $j-1,8,6,6$, all at least
two. The number of reflections is $(6-j)+3=A$.
For $A=7$, write $D'=DR_1^{-1}$ and take
\[
 (P_3,t_3,t_4,t_5,R_7,R_8,R_9,
                        C^{-1},D',R_1,D^{-1}).
 \tag{F1I05.9}
\]
Here $D'$ still has rank five, so there are exactly four
nonreflection entries. The original matrices $C$ and $D$
are products of the displayed factors in every tuple.
Consequently the generated group contains the already proved
$\langle C,D\rangle=G(3,1,9)$, and it is contained in
that same group. This is an actual recovery of generators.

The permutation indices of the first block sum to
$(j-1)+(6-j)=5$; the other three nonreflection permutations
have indices five, three and three. Diagonal reflections
have index zero. The total permutation index is therefore
sixteen in all four tuples. Riemann--Hurwitz on the actual
degree-nine quotient gives $g(Y)=1-9+16/2=0$.
Moving each removed fixed-letter color into its displayed
diagonal reflection preserves both primitive fractional-color
sums. The pure cycle splittings introduce no color.
Both sums thus remain six, as in I01, and the two actual
holomorphic multiplicities are $-1+6=5$.
The four reduced uncolored-cycle lists are
\[
\begin{array}{c|c|c|c}
 A&\hbox{counts over the four nonreflection values}&
       \deg Q&\operatorname{rank}d\operatorname{Per}\\ \hline
 4&(5,1,3,3)&12&5\\
 5&(6,1,3,3)&13&6\\
 6&(7,1,3,3)&14&7\\
 7&(7,1,4,3)&15&8
\end{array}
 \tag{F1I05.10}
\]
The first count is at least $A$ in every row, so paragraph3
places every smooth point of each such ramification profile
in the first rank row of (F1I05.2). Riemann existence for
these product-one generating tuples supplies connected
covers and their faithful monomial quotient towers. Thus
the four rows are nonempty, and their image dimensions equal
the indicated constant differential ranks.

\paragraph{5. Exact remaining locus and comparison maps.}
For $A=4$, a point in the second rank row must have the
uncolored list $(3,3,3,3)$: each count is at most three
and their sum is twelve. Its coprime $\alpha,\beta$ have
degree three. Indeed the I03 total count gives
$12\le c_0+4(3-c_0)=12-3c_0$, so $c_0=0$.
The actual degree-three map $\phi$ consequently satisfies
\[
                  Q=\phi^*B
\]
as reduced divisors: all twelve distinct points lie in
the degree-twelve pullback. This proves a necessary
factorization of the marked-divisor map, with the original
degree-nine $h$ retained. It asserts no factorization of $h$.

For $A=8$, the count criterion alone yields no profile
from paragraph3. A degree-nine nonreflection has at most
seven uncolored cycles: eight uncolored cycles use at least
eight letters, and the remaining possibilities have moving
rank at most one. The explicit matrix test (F1I05.2) still
computes the exact original rank, eight or nine, at every
given smooth point. Existence and component distribution
of the second row remain unresolved at all retained values
where it has not been excluded.

The literature comparison retains the actual marked-Hurwitz
map and quadratic trace of I20--I21. Filom,
\emph{On the Rank of the Thurston Pullback Map}, IMRN2023,
no.24,21838--21873, DOI10.1093/imrn/rnac305,
Theorem1.1 already gives the lower bound $A$ through that
same trace. Each diagonal reflection contributes one regular
target value with exactly one marked colored preimage; each
transposition reflection contributes one simple critical
point with no marked colored point in its fibre. These are
exactly its two counted sets, and the marked source has
twelve points, so its bound is $\min(A,9)=A$ here.
Together with source dimension $A+1$, that published bound
already gives the two possible integers. Its full proof and
Lemma3.1 were read for I04; their exact source scope is
retained. The additional assertions proved here are the
actual polynomial test selecting the integer and the four
explicit nonempty profiles attaining the larger rank.
Fresh Hexagon, arXiv and journal searches are recorded in
the source audit, without a priority assertion.

All operations occur on the original finite tower and
its prescribed marked divisor. The evaluation sequence,
unit quotient and polynomial multiplication give the actual
F02 tangent and Hilbert comparison maps. The original
polarized $q^*,q_*$ comparison retains the F03 theta
multiplier, and the whole period image is the original
F01 variational image. Their exact vector-bundle sequences
give the F04 $K_0$ identities, with Adams operations and
the same cycle map. The four newly evaluated nonempty
profiles narrow T4-INDIVIDUAL; no detector dominance or
additional family outside the original tower is asserted.
